\chapter{Introduction}
\label{ch:introduction}
Much rail maintenance still runs on fixed intervals: a component is inspected or replaced after
a set mileage or time, regardless of how hard it has actually been used. That wastes healthy
parts and still misses parts that fail early under unusual load. Predictive maintenance ties
each intervention to the measured condition of the equipment instead, promising less wasted
material and fewer unnecessary replacements. The friction brakes of a metro train are a good
place to try this: the train brakes at almost every stop, hundreds of times a day, so the brake
components are cycled constantly and any degradation should leave a large, repeating signal.
This report analyses the MetroAT dataset~\cite{MetroAT:2026}: one year of 1\,Hz recordings from a
single Wiener Linien metro train, covering the pneumatic system together with train speed and mode
signals, and annotated with documented failure and maintenance events. All sensor values are
min-max normalised, so the analysis works in relative rather than physical units. Following the
project proposal, the work runs in two phases along the CRISP-DM cycle~\cite{WirthHipp:00}. The first recognises the
train's operational regimes without supervision, clustering the motion signal into a small set of
interpretable states. The second examines the braking that dominates mechanical wear: individual
deceleration events are extracted, their intensity characterised, and braking behaviour tested for
any trend ahead of documented brake-system failures.

Four questions run through the report: (1) whether distinct operational regimes emerge from the
motion data alone, without fixing their number in advance; (2) given those regimes, whether the
air-system sensors can recover the deceleration regime using only signals that did not help define
it; (3) whether deceleration intensity falls into discrete classes or forms a continuum, and how
much of it the air-system sensors can explain; and (4) whether braking behaviour carries any
detectable signal ahead of documented brake-system failures. The remainder of the report is organised by topic rather than by project phase.

\textbf{Structure.} Chapter~\ref{ch:preprocessing} describes how the raw recordings become the
analysis units and how the sensors are partitioned into feature roles.
Chapter~\ref{ch:states} derives and validates four motion-only operational states.
Chapter~\ref{ch:classification} tests whether the air-system sensors alone can recover the
deceleration state. Chapter~\ref{ch:braking} extracts individual deceleration events and regresses
their intensity and dissipated energy from the air sensors. Chapter~\ref{ch:failure} examines the
link to documented failures and the maintenance calendar and reports a near-null pre-failure
result. Chapter~\ref{ch:conclusion} draws the findings together and points to what is worth doing
next.

\textbf{Scope.} The conclusions are bounded by the data: a single train, an eight-and-a-half-month
training window, and roughly ten documented failures, eight of them brake-system failures; the
study is exploratory and not meant as a deployable failure detector. Differences between individual
wagons are not analysed: the many near-duplicate per-wagon channels are instead collapsed by
principal component analysis during preprocessing (Chapter~\ref{ch:preprocessing},
Figure~\ref{fig:pca}).